\title{Group Sequence Policy Optimization}

\begin{abstract}
We present Group Sequence Policy Optimization (GSPO), a reinforcement learning algorithm designed for training large language models with greater stability, efficiency, and effectiveness. Distinct from existing approaches that employ token-level importance ratios, GSPO utilizes a sequence-level importance ratio, implementing clipping, rewarding, and optimization over entire sequences. Our results indicate that GSPO not only enhances training efficiency and outperforms the GRPO algorithm, but also brings pronounced stability to Mixture-of-Experts (MoE) reinforcement learning. Additionally, GSPO shows promise for streamlining reinforcement learning system design. These advantages have played a significant role in the substantial performance gains observed in recent Qwen3 model developments.
\end{abstract}

\section{Introduction}
\label{sec:intro}

Reinforcement learning (RL) has become a critical framework for advancing the scalability of language models \citep{o1,dpsk-r1,qwq32b,qwen3}. Leveraging RL at a large scale empowers language models to address intricate tasks, including high-level mathematics and programming, by enabling more complex and sustained reasoning.

Achieving successful RL scaling in tandem with increased computational resources requires, above all else, training processes that are both stable and resilient. Despite this necessity, current leading RL methods—such as GRPO \citep{grpo}—suffer from pronounced instability when applied to extensive language model training, frequently causing dramatic and irrevocable collapse of models \citep{qwen3, minimax-m1}. Such instability significantly obstructs efforts to expand the capabilities of language models through further RL-driven training.

This study reveals that the instability observed in GRPO originates from the improper application and breakdown of importance sampling weights within its algorithmic framework. This misstep results in significant training variance, which intensifies as the generated response length increases and is further exacerbated by the clipping process, culminating in eventual model failure.

To overcome these fundamental challenges, we introduce \textbf{Group Sequence Policy Optimization (GSPO)}, a novel reinforcement learning algorithm tailored for large language model training. GSPO’s principal innovation is the establishment of a theoretically principled importance ratio, founded on sequence likelihood \citep{click}, which adheres to the core tenets of importance sampling. Furthermore, GSPO determines normalized rewards as the advantage values across multiple responses to a single query, guaranteeing proper sequence-level reward attribution and optimization consistency.

Through experimental analysis, we establish that GSPO markedly outperforms GRPO in terms of stability during training, operational efficiency, and overall effectiveness. Notably, GSPO intrinsically addresses the instability often encountered in RL training for expansive Mixture-of-Experts (MoE) architectures, thus obviating the necessity for intricate stabilization measures and offering prospects for streamlining RL frameworks. These strengths of GSPO have directly facilitated the remarkable gains observed in the newest Qwen3 models. We anticipate that GSPO will serve as a dependable and scalable cornerstone for further developments in large-scale RL training with language models.

\section{Preliminaries}

\paragraph{Notation}
Throughout this work, we represent an autoregressive language model with parameters $\theta$ as the policy $\pi_\theta$. The variable $x$ is used for a query, while $\mathcal{D}$ refers to the collection of queries. For a given response $y$ corresponding to query $x$, its probability according to $\pi_\theta$ is expressed as $\pi_\theta (y | x)=\prod_{t=1}^{|y|} \pi_\theta (y_t | x, y_{<t} )$, where $|y|$ counts the tokens in $y$. Each query-response pair $(x, y)$ is evaluated by a verifier $r$, producing a reward $r(x, y)$ that lies within the interval $[ 0, 1 ]$.

\paragraph{Proximal Policy Optimization (PPO)} Using samples generated from the old policy $\pi_{\theta_\text{old}}$, PPO \citep{ppo} constrains the policy update within a proximal region of the old policy through the clipping mechanism. Specifically, PPO employs the following objective for policy optimization (we omit the KL regularization term hereinafter for brevity, as it is not the focus of this paper):

\begin{align}
\mathcal{J}_\text{PPO}(\theta) = \mathbb{E}_{ x \sim \mathcal{D},\, y \sim \pi_{\theta_\text{old}}( \cdot | x) }
\left[ \frac{1}{|y|} \sum_{t=1}^{|y|} 
\min \left( w_{t}(\theta) \widehat{A}_{t},  \, \mathrm{clip} \left( w_{t}(\theta), 1 - {\varepsilon}, 1 + {\varepsilon}\right) \widehat{A}_{t} \right)
\right],
\end{align}

The importance ratio for the token $y_t$ is given by $
w_{t}(\theta) = \frac{ \pi_{\theta} (y_{t} | x, y_{<t}) }{ \pi_{\theta_\text{old}} (y_{t} | x,y_{<t})} $, where the advantage $\widehat{A}_{t}$ associated with $y_t$ is computed using a separate value model, and $\varepsilon$ denotes the threshold used to clip the importance ratios.

A primary difficulty encountered with PPO in real-world applications stems from its substantial dependency on the value model. Typically, the value model possesses a scale comparable to that of the policy model, resulting in significant demands for memory and computation. Moreover, the algorithm’s performance is critically dependent on the quality of the estimated values. Obtaining dependable value estimations is a nontrivial task, and making the value model robust enough for longer outputs and complex tasks proves to be even more problematic.

\paragraph{Group Relative Policy Optimization (GRPO)} GRPO \citep{grpo} bypasses the need for the value model by computing the relative advantage of each response within a group of responses to the same query. Specifically, GRPO optimizes the following objective:

\begin{align}
\label{equ:grpo}
\mathcal{J}_\text{GRPO}(\theta) = \mathbb{E}_{ x \sim \mathcal{D},\, \{y_i\}_{i=1}^G \sim \pi_{\theta_\text{old}}( \cdot | x) }
\left[ \frac{1}{G} \sum_{i=1}^{G} \frac{1}{|y_i|} \sum_{t=1}^{|y_i|} 
\min \left( w_{i,t}(\theta) \widehat{A}_{i,t},  \, \mathrm{clip} \left( w_{i,t}(\theta), 1 - {\varepsilon}, 1 + {\varepsilon}\right) \widehat{A}_{i,t} \right)
\right],
\end{align}

where $G$ is the number of generated responses to each query $x$ (i.e., the group size), and the importance ratio $w_{i,t}(\theta)$ and advantage $\widehat{A}_{i,t}$ of token $y_{i,t}$ are:

\begin{align}
    w_{i,t}(\theta)=\frac{ \pi_{\theta} (y_{i,t} | x, y_{i,<t}) }{ \pi_{\theta_\text{old}} (y_{i,t} | x,y_{i,<t})},\quad\ 
    \widehat{A}_{i,t} = \widehat{A}_{i} = \frac{r(x, y_i) - \mathrm{mean} \left( \{ r(x, y_i) \}_{i=1}^G \right) }{ \mathrm{std} \left( \{ r(x, y_i) \}_{i=1}^G \right) },
\end{align}

respectively, where all the tokens in $y_i$ share the same advantage as $\widehat{A}_{i}$.

\section{Motivation}
\label{sec:motivation}

As models continue to increase in size, incorporate greater sparsity (such as in Mixture-of-Experts architectures), and produce longer responses, reinforcement learning demands larger rollout batch sizes to fully exploit available hardware resources. To enhance the efficiency with which training samples are used, it is common to subdivide these extensive batches of rollout data into several mini-batches for the purpose of performing gradient updates. However, this introduces an off-policy learning scenario; here, the outputs $y$ originate from a previous policy $\pi_{\theta_\text{old}}$ rather than the updated policy $\pi_\theta$ being refined. This scenario underscores the importance of the clipping strategy employed in algorithms like PPO and GRPO, which serves to restrict the influence of samples that deviate significantly from the current policy in the estimation of gradients.

While mechanisms like clipping aim to manage this off-policy discrepancy, we identify a more fundamental issue in GRPO: \textit{its objective is ill-posed}. This problem becomes particularly acute when training large models on long-response tasks, leading to catastrophic model collapse. The ill-posed nature of the GRPO objective stems from a misapplication of importance sampling weights. The principle of importance sampling is to estimate the expectation of a function $f$ under a target distribution $\pi_\text{tar}$ by re-weighting samples drawn from a behavior distribution $\pi_\text{beh}$:

\begin{align}
\mathbb{E}_{z \sim \pi_\text{tar}} \left[ f(z) \right] = \mathbb{E}_{z \sim \pi_\text{beh}} \left[ \frac{ \pi_\text{tar} (z) }{ \pi_\text{beh} (z)} \, f(z) \right].
\end{align}

Importantly, effective correction for the mismatch between the behavior distribution $\pi_\text{beh}$ and the target distribution $\pi_\text{tar}$ via the importance weight $\frac{ \pi_\text{tar} (z) }{ \pi_\text{beh} (z)}$ necessitates averaging across a substantial number of samples ($N \gg 1$) drawn from $\pi_\text{beh}$. 

By comparison, GRPO introduces the importance weight $\frac{ \pi_{\theta} (y_{i,t} | x, y_{i,<t}) }{ \pi_{\theta_\text{old}} (y_{i,t} | x,y_{i,<t})}$ at each token position $t$, relying solely on a single sampled token $y_{i,t}$ from the next-token distribution $\pi_{\theta_\text{old}}( \cdot | x, y_{i, <t})$. This approach falls short of providing meaningful distributional correction and instead injects considerable variance into the gradient estimates during training. Over lengthy sequences, the variance accumulates, further amplified by the clipping operation. Empirically, this phenomenon precipitates catastrophic model collapse—a failure state from which recovery is impossible, regardless of interventions such as reverting to earlier checkpoints, meticulous adjustment of hyperparameters (including clipping thresholds), increasing sequence length, or changing RL queries.

These empirical findings underscore an intrinsic flaw in GRPO. The inability of token-level importance weighting to rectify the distributional difference highlights a crucial principle: \textbf{the optimization objective must be aligned with the reward's granularity}. Because rewards are distributed at the sequence level, token-based correction mechanisms are insufficient and potentially harmful. Consequently, we are motivated to move away from token-level objectives, instead adopting importance weighting and direct optimization at the \textit{sequence level}.

\section{Algorithm}

\subsection{GSPO: Group Sequence Policy Optimization}

Although the token-level importance weight $\frac{ \pi_{\theta} (y_{i,t} | x, y_{i,<t}) }{ \pi_{\theta_\text{old}} (y_{i,t} | x,y_{i,<t})}$ introduces challenges in GRPO, it is notable that, for language generation tasks, the \textit{sequence-level} importance weight $\frac{ \pi_{\theta} (y | x) }{ \pi_{\theta_\text{old}} (y | x)}$ possesses a well-defined theoretical interpretation. Specifically, it quantifies the divergence of a response $y$ drawn from $\pi_{\theta_\text{old}} (\cdot | x)$ relative to $\pi_{\theta} (\cdot | x)$, providing an intrinsic correspondence to the sequence-level reward signal and functioning as an informative measure for the application of the clipping procedure.

Based on this straightforward observation, we propose the \textbf{Group Sequence Policy Optimization (GSPO)} algorithm. GSPO employs the following sequence-level optimization objective:

\begin{align}
\mathcal{J}_\text{GSPO} (\theta) =
\mathbb{E}_{ x \sim \mathcal{D},\, \{y_i\}_{i=1}^G \sim \pi_{\theta_\text{old}}( \cdot | x) }
\left[ 
\frac{1}{G} \sum_{i=1}^{G}
\min \left( s_{i}(\theta)  \widehat{A}_{i},  \, \mathrm{clip} \left( s_{i}(\theta), 1 - {\varepsilon}, 1 + {\varepsilon} \right) \widehat{A}_{i} \right) 
\right],
\label{equ:gspo}
\end{align}

where we adopt the group-based advantage estimation:

\begin{align}
\widehat{A}_{i} = \frac{r(x, y_i) - \mathrm{mean} \left( \{ r(x, y_i) \}_{i=1}^G \right) }{ \mathrm{std} \left( \{ r(x, y_i) \}_{i=1}^G \right) },
\end{align}

and define the importance ratio $s_{i}(\theta)$ based on sequence likelihood \citep{click}:

\begin{align}
s_{i}(\theta) = \left( \frac{ \pi_{\theta} (y_i | x) }{ \pi_{\theta_\text{old}} (y_i | x)} \right)^{\frac{1}{|y_i|}}
=
\exp \left( \frac{1}{|y_i|} \sum_{t=1}^{|y_i|} \log \frac{ \pi_{\theta} (y_{i,t} | x, y_{i,<t}) }{ \pi_{\theta_\text{old}} (y_{i,t} | x,y_{i,<t})} \right).
\end{align}

As a result, GSPO implements clipping at the response level rather than at the token level, thereby filtering out excessively ``off-policy'' samples from the gradient calculation. This approach aligns with the reward mechanism and optimization that operate at the sequence scale. To address variance and maintain $s_{i}(\theta)$ within a consistent numerical scope, we introduce length normalization in $s_{i}(\theta)$. Without this normalization, the probability shifts for just a few tokens could lead to substantial changes in the sequence-level importance ratio, causing responses of various lengths to necessitate different clipping thresholds. Additionally, we observe that the clipping intervals used in GSPO and those found in earlier methods (such as GRPO) often differ by orders of magnitude, which is a consequence of their differing definitions for importance ratios.

\subsection{Gradient Analysis}
\label{subsec:gradient}

We can derive the gradient of the GSPO objective as follows (clipping is omitted for brevity):

\begin{align}
\nabla_{\theta} \mathcal{J}_\text{GSPO} (\theta)
=&\ 
\nabla_{\theta} \mathbb{E}_{ x \sim \mathcal{D},\, \{y_i\}_{i=1}^G \sim \pi_{\theta_\text{old}}( \cdot | x) }
\left[ 
\frac{1}{G} \sum_{i=1}^{G} s_{i}(\theta) \widehat{A}_{i}
\right] \\
= &\ 
\mathbb{E}_{ x \sim \mathcal{D},\, \{y_i\}_{i=1}^G \sim \pi_{\theta_\text{old}}( \cdot | x) }
\left[ 
\frac{1}{G} \sum_{i=1}^{G}
s_{i}(\theta) \widehat{A}_{i}
\cdot \nabla_{\theta} \log s_{i}(\theta)
\right] \\
=&\ 
\mathbb{E}_{ x \sim \mathcal{D},\, \{y_i\}_{i=1}^G \sim \pi_{\theta_\text{old}}( \cdot | x) }
\left[ 
\frac{1}{G} \sum_{i=1}^{G} 
\left( \frac{ \pi_{\theta} (y_i | x) }{ \pi_{\theta_\text{old}} (y_i | x)} \right)^{\frac{1}{|y_i|}} \widehat{A}_{i} 
\cdot \frac{1}{|y_i|} \sum_{t=1}^{|y_i|} \nabla_{\theta} \log \pi_{\theta} (y_{i,t} | x, y_{i,<t}) 
\right].
\label{equ:gspo_grad}
\end{align}

For comparison, the gradient of the GRPO objective is as follows (note that $\widehat{A}_{i,t} = \widehat{A}_{i}$):

\begin{align}
\nabla_{\theta} \mathcal{J}_\text{GRPO} (\theta) 
=&\ 
\nabla_{\theta} \mathbb{E}_{ x \sim \mathcal{D},\, \{y_i\}_{i=1}^G \sim \pi_{\theta_\text{old}}( \cdot | x) }
\left[ \frac{1}{G} \sum_{i=1}^{G} \frac{1}{|y_i|} \sum_{t=1}^{|y_i|} 
w_{i,t}(\theta) \widehat{A}_{i,t}
\right] \\
=&\
\mathbb{E}_{ x \sim \mathcal{D},\, \{y_i\}_{i=1}^G \sim \pi_{\theta_\text{old}}( \cdot | x) }
\left[ \frac{1}{G} \sum_{i=1}^{G} \widehat{A}_{i} 
\cdot \frac{1}{|y_i|} \sum_{t=1}^{|y_i|} 
\frac{ \pi_{\theta} (y_{i,t} | x, y_{i,<t}) }{ \pi_{\theta_\text{old}} (y_{i,t} | x,y_{i,<t})} 
\nabla_{\theta} \log \pi_{\theta} (y_{i,t} | x, y_{i,<t})  
\right].
\end{align}

As a result, the primary difference between GSPO and GRPO is rooted in \textit{the method used to assign weights to the gradients of token log likelihoods}. GRPO applies individualized weighting to each token based on its corresponding ``importance weight'' $\frac{ \pi_{\theta} (y_{i,t} | x, y_{i,<t}) }{ \pi_{\theta_\text{old}} (y_{i,t} | x,y_{i,<t})}$. These non-uniform weights, which may span the ranges $(0, 1+\varepsilon]$ for $\widehat{A}_{i} >0$, or $[1-\varepsilon, +\infty)$ for $\widehat{A}_{i} < 0$, introduce nontrivial variability. This variability can build up over training, potentially resulting in unintended outcomes. Conversely, GSPO assigns identical weights to every token within a generated response, thereby removing the source of instability inherent to GRPO.

\subsection{GSPO-token: A Token-level Objective Variant}

In scenarios like multi-turn RL, we may desire a finer-grained advantage adjustment than the sequence level. To this end, we introduce a token-level objective variant of GSPO, namely \textbf{GSPO-token}, to allow token-wise advantage customization:

\begin{align}
\mathcal{J}_\text{GSPO-token}(\theta)
=
\mathbb{E}_{ x \sim \mathcal{D},\, \{y_i\}_{i=1}^G \sim \pi_{\theta_\text{old}}( \cdot | x) }
\left[ 
\frac{1}{G} \sum_{i=1}^{G} \frac{1}{|y_i|} \sum_{t=1}^{|y_i|} 
\min \left( s_{i,t}(\theta) \widehat{A}_{i,t},  \, \mathrm{clip} \left( s_{i,t}(\theta), 1 - {\varepsilon}, 1 + {\varepsilon} \right) \widehat{A}_{i,t} \right) 
\right],
\label{equ:gspo-token}
\end{align}

where

\begin{align}
s_{i,t}(\theta) 
= \mathrm{sg} \left[ s_{i}(\theta) \right]  \cdot \frac{ \pi_{\theta} (y_{i,t} | x, y_{i,<t}) }{ \mathrm{sg} \left[ \pi_{\theta} (y_{i,t} | x, y_{i,<t}) \right] },
\end{align}

and $\mathrm{sg}[\cdot]$ denotes only taking the numerical value but stopping the gradient, corresponding to the \texttt{detach} operation in PyTorch. The gradient of GSPO-token can be derived as:

\begin{align}
\nabla_{\theta} \mathcal{J}_\text{GSPO-token}(\theta)
=&\ 
\nabla_{\theta} \mathbb{E}_{ x \sim \mathcal{D},\, \{y_i\}_{i=1}^G \sim \pi_{\theta_\text{old}}( \cdot | x) }
\left[ 
\frac{1}{G} \sum_{i=1}^{G}
\frac{1}{|y_i|} \sum_{t=1}^{|y_i|} 
s_{i,t}(\theta) \widehat{A}_{i,t}
\right] \\
=&\ 
\mathbb{E}_{ x \sim \mathcal{D},\, \{y_i\}_{i=1}^G \sim \pi_{\theta_\text{old}}( \cdot | x) }
\left[ 
\frac{1}{G} \sum_{i=1}^{G} s_i (\theta) \cdot \frac{1}{|y_i|} \sum_{t=1}^{|y_i|} 
\widehat{A}_{i,t} \frac{ \nabla_{\theta} \pi_{\theta} (y_{i,t} | x, y_{i,<t}) }{ \pi_{\theta} (y_{i,t} | x, y_{i,<t}) }
\right] \\
=&\ 
\mathbb{E}_{ x \sim \mathcal{D},\, \{y_i\}_{i=1}^G \sim \pi_{\theta_\text{old}}( \cdot | x) }
\left[ 
\frac{1}{G} \sum_{i=1}^{G}  \left( \frac{ \pi_{\theta} (y_i | x) }{ \pi_{\theta_\text{old}} (y_i | x)} \right)^{\frac{1}{|y_i|}}
\cdot \frac{1}{|y_i|} \sum_{t=1}^{|y_i|}
\widehat{A}_{i,t} \nabla_{\theta} \log \pi_{\theta} (y_{i,t} | x, y_{i,<t})  
\right].
\label{equ:gspo-token_grad}
\end{align}

Observe that the expression $\frac{ \pi_{\theta} (y_{i,t} | x, y_{i,<t}) }{ \mathrm{sg} \left[ \pi_{\theta} (y_{i,t} | x, y_{i,<t}) \right] }$ evaluates to 1, implying that $s_{i,t}(\theta)$ is numerically equivalent to $s_{i}(\theta)$. By analyzing Equation~\eqref{equ:gspo} alongside~\eqref{equ:gspo-token}, and Equation~\eqref{equ:gspo_grad} alongside~\eqref{equ:gspo-token_grad}, it becomes clear that GSPO-token and GSPO produce identical numerical results in terms of the optimization objective, clipping mechanism, and theoretical gradient provided the advantages for all tokens in the output $y_i$ are set uniformly (specifically, $\widehat{A}_{i,t} = \widehat{A}_{i}$). Nonetheless, GSPO-token offers enhanced adaptability by permitting distinct advantage values for each individual token.

\section{Experiments and Discussion}

\subsection{Empirical Results}

We utilize a cold-start model fine-tuned from Qwen3-30B-A3B-Base for our experiments, presenting both the training reward trajectories and performance metrics on the AIME'24 (mean Pass@1 calculated from 32 samples), LiveCodeBench (202410-202502, mean Pass@1 from 8 samples), and CodeForces (Elo Rating) evaluation suites. In our RL training setup, each rollout batch is subdivided into four mini-batches to perform gradient updates. For GSPO, the left and right clipping parameters in Equation~\eqref{equ:gspo} are fixed at 3e-4 and 4e-4, respectively. Our baseline employs GRPO, where we adjust the left and right clipping limits in Equation~\eqref{equ:grpo} to 0.2 and 0.27, respectively; these values are selected with care to guarantee an equitable comparison. It is important to note that GRPO requires the Routing Replay technique to facilitate standard convergence in MoE RL—a detail that will be elaborated upon in \S~\ref{subsec:routing_replay}—whereas, \textit{GSPO eliminates the necessity for this approach}.

As depicted in Figure~\ref{fig:results}, training with GSPO maintains stable progress across the entire process. Our findings indicate that \textbf{GSPO consistently enhances model performance by increasing training compute, periodically refreshing the query set, and lengthening generated sequences}. Furthermore, GSPO outperforms GRPO in terms of training efficiency, yielding superior training accuracy and benchmark results for equivalent training compute and query usage. Additionally, GSPO has been successfully utilized in RL training for the newest Qwen3 models, providing robust evidence of its effectiveness in maximizing RL scaling benefits for large language models.

\subsection{Curious Observation on Clipping Fractions}

GSPO sets itself apart from GRPO by implementing clipping at the level of entire sequences, rather than applying it to individual tokens. As illustrated in Figure~\ref{fig:clipping}, this leads to a marked contrast: the proportion of tokens clipped in GSPO is roughly two orders of magnitude higher than in GRPO, a gap that persists even when clipping ranges are modified. Notably, GSPO maintains greater training efficiency compared to GRPO, despite the fact that it clips far more tokens and, as a result, utilizes a smaller subset for training or gradient calculation. This seemingly paradoxical result—that extensive token clipping in GSPO correlates with improved training outcomes—suggests that GRPO’s token-wise gradient estimates are prone to excessive noise and lack efficacy in sample utilization. On the other hand, GSPO's strategy of leveraging gradients at the sequence level yields a more robust and actionable learning signal.

\subsection{Benefit of GSPO for MoE Training}
\label{subsec:routing_replay}

\paragraph{Background}
Relative to reinforcement learning (RL) applied to dense models, MoE models pose distinct challenges regarding training stability due to their sparsely activated architecture. Notably, our experiments reveal that employing the GRPO algorithm leads to \textit{expert-activation volatility} in MoE models, which may impede proper convergence during RL training. Specifically, following one or several updates of the model's parameters, the subset of experts activated in response to the same input may shift appreciably. For instance, with the Qwen3-30B-A3B-Base model comprising 48 layers, each RL gradient update yields around 10\% of experts—considering any given rollout sample—whose activations differ between the updated policy $\pi_{\theta}$ and the preceding policy $\pi_{\theta_\text{old}}$.

This kind of fluctuation, more acute in deeper MoE architectures, leads to high variability in the token-level importance weights $w_{i,t}(\theta) = \frac{ \pi_{\theta} (y_{i,t} | x, y_{i,<t}) }{ \pi_{\theta_\text{old}} (y_{i,t} | x,y_{i,<t})}$. Such instability undermines the utility of these ratios, as elaborated in \S~\ref{sec:motivation} and~\ref{subsec:gradient}, and ultimately disrupts the standard convergence patterns expected in RL training.

\paragraph{Our Previous Approach} In earlier work, we addressed this problem using the \textbf{Routing Replay} method during training. This approach involves storing the experts selected by $\pi_{\theta_\text{old}}$ and then replicating these same routing configurations in $\pi_{\theta}$ when calculating the importance ratios $w_{i,t}(\theta) = \frac{ \pi_{\theta} (y_{i,t} | x, y_{i,<t}) }{ \pi_{\theta_\text{old}} (y_{i,t} | x,y_{i,<t})}$. As a result, both $\pi_{\theta} (y_{i,t} | x, y_{i,<t})$ and $\pi_{\theta_\text{old}} (y_{i,t} | x,y_{i,<t})$ utilize the identical activated expert networks for every token $y_{i,t}$. This alignment contributes to stabilizing the token-level importance ratio and facilitates the consistent optimization of the same expert set over consecutive gradient steps. As depicted in Figure~\ref{fig:routing_replay}, Routing Replay proves to be a crucial component for achieving standard convergence within GRPO training applied to MoE architectures.

\paragraph{Benefit of GSPO} While Routing Replay allows GRPO training of MoE models to reach convergence, it does so at the expense of increased memory usage and communication costs, and also restricts the MoE's accessible capacity by repeatedly using prior routing patterns. In comparison, GSPO, as illustrated in Figure~\ref{fig:results}, operates independently of Routing Replay. GSPO accurately computes the importance ratios $s_i(\theta)$ using standard procedures, ensuring typical convergence and stable optimization dynamics. The principal observation is that GSPO relies exclusively on the sequence likelihood (i.e., $\pi_{\theta} (y_i | x)$) and exhibits robustness to variations in per-token likelihoods (i.e., $\pi_{\theta} (y_{i,t} | x, y_{i,<t})$). Because the MoE model preserves its language modeling performance throughout, the sequence likelihood remains relatively stable. Ultimately, GSPO addresses the challenge of expert activation instability inherent to MoE models, eliminating the necessity for auxiliary mechanisms such as Routing Replay. This streamlines the training workflow, enhances stability, and enables unrestricted utilization of the model's capacity.

\subsection{Benefit of GSPO for RL Infrastructure}

Due to the differences in precision between training frameworks such as Megatron and inference systems like SGLang and vLLM, it is common practice to recalculate the likelihoods of generated responses under the previous policy $\pi_{\theta_\text{old}}$ using the training framework. Nonetheless, GSPO employs optimization based on sequence-level likelihoods instead of relying on token-level likelihoods. Since sequence-level likelihoods exhibit greater robustness to minor precision variations, GSPO enables direct utilization of likelihood values from the inference engine for optimization purposes. This approach eliminates the requirement for recomputation using the training engine. As a result, GSPO proves particularly advantageous in contexts such as partial rollout, multi-turn RL, and settings with separated training and inference processes.

\section{Conclusion}

We introduce Group Sequence Policy Optimization (GSPO), a novel reinforcement learning algorithm designed to enhance the training of large language models. GSPO adheres to the core concepts of importance sampling by calculating importance ratios using sequence likelihood, and subsequently employs sequence-level clipping, reward assignment, and optimization. Compared to GRPO, GSPO achieves substantially greater training stability, efficiency, and overall performance, and demonstrates significant advantages in large-scale RL applications, particularly for MoE model training. This algorithm has been instrumental in driving the recent advancements observed in the Qwen3 models. As GSPO provides a scalable basis for algorithmic development, our ongoing efforts focus on further scaling reinforcement learning, anticipating that this will lead to major breakthroughs in artificial intelligence.

\end{document}